% Appendix B. Deferred proof details -- 950 words
% (appendix share 12.3% of 14,250; option (a), O-6)
%
% B.1 Serialization of the P-hardness reduction (C1.4): record layout,
%     field widths, preliminary log-space validation, and the map from
%     malformed inputs to the fixed feasible negative instance ... 180 words
% B.2 Separation estimates (C7): radius, determinant estimates and M
%     .......................................................... 120 words
% B.3 Zero probes (C10-C13; C13 with its compact-closure and rate-fiber
%     arguments) ............................................... 400 words
% B.4 One-state compilers and bounds (C3 compiler and polynomial-size
%     proof; C5 template and encoding bounds) ................... 250 words

\section{Deferred proof details}\label{app:proofs}

\subsection{Proofs for Section~\ref{sec:problem}}\label{app:problem-proofs}

\begin{proof}[Proof of~\eqref{eq:word-joint}]
For each fixed base trajectory, summing the register factors backwards uses
their unit row sums. The remaining sum over base trajectories is the
indicated base-matrix product. The base kernels share the uniform and probe
projections, which gives the displayed expansion.
\end{proof}

\begin{proof}[Proof of Lemma~\ref{lem:compile}]
By~\eqref{eq:word-joint}, each side combines the pairwise orthogonal
nonzero functions $1$ and $\chi_i(g)\chi_i(h)$ of $(g,h)$. Compare
coefficients, and in~(a) divide by $\theta(w)=\theta(w')>0$.
\end{proof}

\subsection{Proofs for Section~\ref{sec:core}}\label{app:core-proofs}

% compressed 2026-11-24: type (c): in part (a) the subdivision sentence is
% named "rational subdivision" and later proofs (Lemma transfer(d), the
% control, the algebraic transfer, Theorem rational-roots) cite it instead
% of repeating it; type (e): parts (b) and (c) tightened, every step kept.
\begin{proof}[Proof of Lemma~\ref{lem:restriction}]
(a) For $t>0$, an active test gains $t^{\eta\cdot a}$ and
$t^{\eta\cdot b}$ on its two sides, equal as $\Delta_S\eta=0$. An inactive
test has both sides zero at~$x$, through a coordinate outside $S$ or
through $c=0$, and they stay zero along the curve, as those coordinates
do. The equalities pass to $t=0$ by continuity, and $|x_j(t)|\le|x_j|$
keeps the budgets. The curve is rational at rational~$t$ and uniformly
continuous on $[0,1]$, so finely spaced rational samples give, for each
$\varepsilon>0$, a sequence from $x$ to $x(0)$ with steps
below~$\varepsilon$ (\emph{rational subdivision}).

(b) A test active on $U$ is active on~$S$, and $a-b$ vanishes outside~$U$,
so $\Delta_U\eta|_U=0$, with zero set $U\cap T$. Along
$S=S_0\to S_1\to\dots\to S_m=T$ this gives
$U\to U\cap S_1\to\dots\to U\cap S_m=U\cap T$.

(c) Let $C$ have least cardinality among the supports $S$ peels to. If $S$
peels to~$U$, then by~(b) $S$ peels to $U\cap C$, so $C\subseteq U$ by
minimality. Thus $C$ is least and unique, and $U$ peels to $U\cap C=C$,
which is also least for~$U$. Supports with the same core both peel to it;
conversely, a common support they both peel to has one core, that of each.
\end{proof}

% compressed 2026-11-24: type (e): the right-inverse step is stated in two
% sentences instead of four; no step removed.
\begin{proof}[Proof of Lemma~\ref{lem:limit}]
By convergence, $T=\supp x^*$ keeps the signs of~$\sigma$. Put
$Z=S\setminus T$; if $Z=\emptyset$, take $\eta=0$.
Otherwise let $\ell^{(k)}=(\log|x^{(k)}_j|)_{j\in S}$. Every test active
on $S$ has $c\neq0$ and nonzero monomials at $x^{(k)}$, so
\[
  \Delta_{S,T}\,\ell^{(k)}_T+\Delta_{S,Z}\,\ell^{(k)}_Z=\gamma ,
\]
where $\Delta_{S,T}$ and $\Delta_{S,Z}$ are the column blocks of $\Delta_S$
and $\gamma$ lists the $\log|c|$. This system includes a row for every test active on~$S$,
including those whose monomials vanish at~$x^*$. As $\ell^{(k)}_T$
converges, $\Delta_{S,Z}\ell^{(k)}_Z$ is bounded, while
$\ell^{(k)}_Z\to-\infty$ entrywise. For a fixed linear right inverse $B$
of $\Delta_{S,Z}$ on its image,
$\ell^{(k)}_Z-B\Delta_{S,Z}\ell^{(k)}_Z\in\ker\Delta_{S,Z}$ is within a
bounded distance of $\ell^{(k)}_Z$, so its negative is eventually
positive. The integer matrix $\Delta_{S,Z}$ has a rational kernel basis,
so rational approximation and clearing denominators give a positive
integer kernel vector. Extending it by zero on~$T$ yields a weight vector
with $\Delta_S\eta=0$ and zero set~$T$. So $S$ peels to~$T$,
$\core(T)=\core(S)$ by Lemma~\ref{lem:restriction}(c), and with the signs
on~$T$, $\SC(x^*)=\SC(x^{(k)})$. The peeling step keeps the starting
magnitudes on the surviving support; those magnitudes may differ from the
limit magnitudes.
\end{proof}

The bridge to Tevelev's criterion~\cite[Lemma~2.2]{tevelev} is as follows.
Impose the tests active on $S$ in the complex torus on~$S$. A component
containing a subsequence has algebraic closure meeting the closed orbit of
$\mathbb A^{S\setminus T}\times(\mathbb C^*)^T$, where $T=\supp x^*$.
Tevelev's Lemma~2.2 gives a valuation vector in the relative interior of
the coordinate cone. Valuation of the active equations puts this vector in
$\ker\Delta_S$, and clearing denominators gives the weights. The retained
signs follow from real convergence.

% compressed 2026-11-24: (d) the opening roadmap sentence ("We first join x
% and x' by a real curve ... exactly") is cut; (c) the inactive-test check
% cites Lemma restriction(a); the closing epsilon/3 step is named
% "rational sampling" for reuse in Lemma transfer(c).
\begin{proof}[Proof of Lemma~\ref{lem:interpolation}]
Let $S$ be the common support. On $S$ put
\[
  \xi_j(t)=\sgn(x_j)\,|x_j|^{1-t}|x'_j|^{t}\qquad(0\le t\le1),
\]
and $\xi_j(t)=0$ off~$S$. An active test reads
$(a-b)\cdot\log|\xi|=\log|c|$ on~$S$, with a sign condition independent
of~$t$; this is affine in $t$ and holds at both ends, hence throughout.
Inactive tests stay zero as in the proof of
Lemma~\ref{lem:restriction}(a). Each $|\xi_j(t)|$ is convex in~$t$,
strictly when $|x_j|\neq|x'_j|$, so a budget containing a changed
coordinate has strict slack for $0<t<1$, and a budget of unchanged
coordinates keeps its endpoint value.

Write each $\rho_j=|x'_j|/|x_j|$ $(j\in S)$ as
$\prod_\pi\pi^{\zeta_{j\pi}}$ with integer exponents over one finite set of
primes. The endpoints satisfy the active tests with the same signs, so
$\rho^{a-b}=1$ for each, and unique factorization gives
$\Delta_S\zeta_\pi=0$ for every~$\pi$, with $\zeta_{j\pi}=0$ wherever
$\rho_j=1$. At an interior time $t$ choose positive rationals $r_\pi$ near
$\pi^t$, and put $x_j\prod_\pi r_\pi^{\zeta_{j\pi}}$ on $S$ and $0$
off~$S$. This point is rational, has the signs of~$x$, keeps the unchanged
coordinates, satisfies the inactive tests as before, and satisfies every
active test exactly, each gaining the factor
$\prod_\pi r_\pi^{(a-b)\cdot\zeta_\pi}=1$ between its sides. It tends to
$\xi(t)$ as $r_\pi\to\pi^t$, so close choices meet the budgets.

Finally, given $\varepsilon>0$, uniform continuity gives
$0=t_0<\dots<t_n=1$ with
$\|\xi(t_i)-\xi(t_{i-1})\|_\infty<\varepsilon/3$; keeping $x$ and $x'$ and
replacing each interior $\xi(t_i)$ by such a rational point within
$\varepsilon/3$ gives steps below~$\varepsilon$ (\emph{rational
sampling}).
\end{proof}

% compressed 2026-11-24: (e) sentences merged; both directions kept.
\begin{proof}[Proof of Lemma~\ref{lem:compact-chain}]
Let $h$ and $h'$ lie in different components of~$Z$. In a compact
Hausdorff space a point's component is the intersection of its clopen
neighborhoods, so a clopen $V$ contains $h$ but not~$h'$. The compact sets
$V$ and $Z\setminus V$ are at a positive distance~$d$, and steps below $d$
never leave~$V$.

Conversely, fix $\varepsilon>0$, let $R\subseteq Y$ be the points reached
from $h$ by sequences in $Y$ with steps below~$\varepsilon$, and let $O$
and $O'$ be the unions of the open $\varepsilon/3$-balls about the points
of $R$ and of $Y\setminus R$. They cover~$Z$, as $Y$ is dense, and are
disjoint, since a common point would put a point of $Y\setminus R$ within
$2\varepsilon/3$ of~$R$, hence in~$R$. The component of $h$ is connected
and meets~$O$, so it lies in~$O$; if it contains $h'$, then $h'\in O$, so
$h'\notin O'$ and $h'\in R$.
\end{proof}

\subsection{Proofs for Section~\ref{sec:tractable}}\label{app:tractable-proofs}

% compressed 2026-11-24: type (e): two sentences merged.
\begin{proof}[Proof of Lemma~\ref{lem:separate}]
Lemma~\ref{lem:compile}(a) splits each held word test into the equation
in~(a) and one probe equation per probe. Each value test constrains
$\theta$ or $\alpha$ alone, and every parameter vector satisfies the
reversal equations and the budgets. As $\theta_m\ge\delta>0$ and
$Q_m(g,h)\ge(1-\kappa)/|X|>0$, an entry of $P_m$ is positive exactly when
that of $L_m$ is.
\end{proof}

% compressed 2026-11-24: type (b): in part (a) the quotient bound is stated
% in one sentence; type (c): in part (c) the closing subdivision (finite
% subdivision at epsilon/3, keep endpoints, approximate interior samples
% within epsilon/3) is replaced by a citation of rational sampling (proof
% of Lemma interpolation), and in part (d) the sampling sentence cites
% rational subdivision.
\begin{proof}[Proof of Lemma~\ref{lem:transfer}]
(a) A row of $P_m$ sums to $\theta_m$, as $Q_m$ and $L_m$ are stochastic.
Since $Q_m\chi_i=\alpha_{mi}\chi_i$,
\[
  \chi_i(g)\sum_{h,r'}P_m((g,r),(h,r'))\,\chi_i(h)=\theta_m\alpha_{mi}.
\]
Tables differing by at most $\tau$ per entry change both row sums,
ordinary and signed, by at most $N\tau$, the weights being signs. As
$\theta_m\ge\delta$ at both and $|\alpha_{mi}|\le\kappa<1$, subtracting the
quotients gives
$|\alpha_{A,mi}-\alpha_{B,mi}|\le N\tau/\delta+|\alpha_{B,mi}|\,N\tau/\delta
\le2N\tau/\delta$.

(b) Each entry is $\theta_mQ_m(g,h)L_m(g,h;r,r')$, a product of three
numbers in $[0,1]$, and $Q_m$ is affine in~$\alpha$. On $[0,1]$,
$|uvw-u'v'w'|\le|u-u'|+|v-v'|+|w-w'|$.

(c) If $\theta=\theta'$, keep it fixed. Otherwise normalize
$\gamma_m(t)=\theta_m^{1-t}(\theta'_m)^t$ by $\sum_m\gamma_m(t)$. Each
equation $\theta(w)=c\,\theta(w')$ with $|w|=|w'|$ holds along $\gamma$, as
at both ends, and survives normalization, both sides having equal degree.
Every $\gamma_m(t)\ge\delta$, and for $0<t<1$ the weighted
arithmetic--geometric mean inequality gives $\sum_m\gamma_m(t)<1$, as some
coordinate changes, so every normalized probability exceeds~$\delta$.
Rational sampling (proof of Lemma~\ref{lem:interpolation}) of the ratios
$\theta'_m/\theta_m$, then normalization, keeps the equations exact, and
the slack above the probability floor absorbs the approximation at interior samples.

(d) Unit row sums, stationarity of~$\nu$, reciprocity and the block-mass
conditions are affine in~$L$, and a convex combination of two nonnegative
tables with the same positive entries has those positive entries. At fixed
$\theta$ and $\alpha$ each table entry is linear in~$L$, so rational
subdivision of the segment gives the sequence.
\end{proof}

% compressed 2026-11-24: type (b): the Cramer bound in part (a) is stated
% in one sentence; type (e): the verifier steps of part (c) are merged into
% two sentences, each check kept.
\begin{proof}[Proof of Theorem~\ref{thm:certificates}]
One LP per round finds all removable coordinates at once:
\[
  \max\sum_{j\in S}\zeta_j\quad\text{subject to}\quad
  \Delta_S\eta=0,\quad \eta\ge0,\quad 0\le\zeta_j\le1,\quad \zeta_j\le\eta_j .
\]
A coordinate is removable exactly when some feasible $\eta$ is positive
there. Summing one such $\eta$ per removable coordinate and scaling
attains their number, the maximum, so every optimum has $\zeta_j=1$
exactly on the removable coordinates and $\eta=0$ exactly on the
others. The LP has at most $2n$ nonnegative variables, entries at most
$D_0$ and right sides $0$ or~$1$, so its polyhedron is pointed. By
Cramer's rule an optimal vertex has one positive integer denominator
which, like its numerators, is at most $(2n)!\,D_0^{2n}$ in absolute
value; clearing it gives integer weights. There are at most $n$ deletion
stages, each removing a coordinate, and empty support stops without an LP.
This proves~(a).

For~(b), if $\eta\ge0$ and $\Delta_C\eta=0$, then
$(\Delta_C^{\mathsf T}\lambda)\cdot\eta=\lambda\cdot\Delta_C\eta=0$ forces
$\eta=0$, so $C$ is irreducible, hence the core by
Lemma~\ref{lem:restriction}(c). Conversely, at an irreducible $C$ no
nonzero real $\eta\ge0$ has $\Delta_C\eta=0$, as rational linear
feasibility would give a rational one, so Gordan's alternative gives
$\Delta_C^{\mathsf T}\lambda>0$. For nonempty $C$, express this
combination, scaled to be at least one in every coordinate, in a basis of
the row space of $\Delta_C$ drawn from the difference vectors of its
active tests. Independence makes the multiplier polyhedron pointed, and
Cramer's rule at a vertex, times the absolute determinant, gives integers
of absolute value at most $n!\,D_0^n$. Empty $C$ uses the empty
certificate. At most $n$ stages of $n$ weights, and at most $n$
multipliers with their indices, give the stated binary length.

For~(c), compile $\Phi_p$ with $|M||I|$ coordinates, exponents at most~$H$,
and one test per reversal equation, held calibration, and pair of held word
test and probe. Compute both signed cores and apply
Theorem~\ref{thm:fixed-profile}. For either verdict, supply a certificate
as in~(b) for each endpoint under the same compiled tests. From the
endpoint's support, the verifier checks at each stage that the weights are
nonnegative integers balancing every test active on the current support,
and passes to exactly the zero-weight coordinates. It then checks the Gordan
inequalities at the terminal support and reads the surviving signs from
the endpoint. It accepts exactly when the two signed cores agree for yes or
differ for no. Checking compares sums of $n$ products of numbers of
polynomial binary length.
\end{proof}

% compressed 2026-11-24: type (e): two sentences merged.
\begin{proof}[Proof of Proposition~\ref{prop:long-sequences}]
On the full support the weights satisfy $\eta_{i+1}=2\eta_i$, so
$\eta_i=2^i$ peels every coordinate. The equations have equal weighted
degree on their two sides for weights $2^i$ on~$x_i$; their ordinary
degrees differ. Both signed cores are zero, and
Theorem~\ref{thm:core} applies, since the budgets bound the feasible set.
Ordered from~$0$, a sequence with steps below $1/4$ has a first point with
$x_0\ge1/4$. There $x_0<1/2$, so its reduced denominator is $q\ge2$; hence
$x_k=x_0^{2^k}$ has reduced denominator $q^{2^k}\ge2^{2^k}$, of binary
length at least $2^k+1$.
\end{proof}

% compressed 2026-11-24: (a) the weight equations, stated in the body's
% proof sketch of Theorem p-hard, are cited instead of re-listed; the
% remaining paragraphs are tightened with every case kept.
\begin{proof}[Proof of Theorem~\ref{thm:p-hard}]
Both endpoints are feasible whatever the circuit: the probability floor
forces the label probabilities, every word has length three, the gate tests avoid~$s$,
and $s$ enters the output test squared. The table support is full, as the
base kernels are positive and the register has one value.

Common letters contribute equally to both sides of a weight equation and
the reversal equations are identities, so active tests give the weight
equations listed in the proof sketch after Theorem~\ref{thm:p-hard}. Let
$K$ be the core of the full support, and $S_0$ the set of~$z$, the false
gates and, if the output is~$0$, the sentinel. We show $K=S_0$.

At~$S_0$, every coordinate other than $z$ has an active test tying it
to~$z$. A constant~$0$ has its own test; a false AND gate, the slot of a
false input; a false OR gate, its test (both inputs are false); and $s$,
the output test. So $S_0$ is irreducible, and Lemma~\ref{lem:irreducible}
gives $S_0\subseteq K$.

Conversely, let $\gamma$ be the first true gate in~$K$, if any. Its tests
are identities or inactive on~$K$, as its true inputs come earlier and lie
outside~$K$, and every other test contains $\mathsf g_\gamma$ as a common
letter or not at all. So the unit weight vector at $g_\gamma$ balances
every test active on~$K$, contradicting the irreducibility of the core
(Lemma~\ref{lem:irreducible}). If the output is~$1$, the output test is
then inactive on~$K$, and the same argument removes~$s$. Hence $K=S_0$.

The endpoints have full support and the same signs except at~$s$, so
their signed cores agree exactly when $s\notin K$, that is, when the
output is~$1$; Theorem~\ref{thm:fixed-profile} gives the equivalence.
The log-space transducer of Appendix~\ref{app:serialization} writes the
instance record by record, reading only the requested gate and evaluating
none.
\end{proof}

% compressed 2026-11-24: (e) sentences merged; every case of the cover
% analysis, the admissibility check and the nesting check are kept.
\begin{proof}[Proof of Lemma~\ref{lem:covers}]
Choose feasible configurations at $p$ and~$q$, with parameter vectors
$(\theta,\alpha,L)$ and $(\theta',\alpha',L')$. %
A decision at a level of~$p$ is constant on the finer blocks of~$q$, so
$L'$ meets the block-mass conditions of~$p$. By Lemma~\ref{lem:compile}
and the definitions of the value tests, passing any held test of either
shape depends only on $\theta$ and~$\alpha$. Since $\theta_mQ_m(g,h)>0$,
the table support is determined by~$L$. Each case below gives an attained
profile strictly above $p$ and at most~$q$, which the cover forces to
equal~$q$.

If the supports differ, $L'$ with $\theta$ and $\alpha$ attains the
partitions and tests of~$p$ with the support of~$q$. If the supports agree
and the held word tests differ, the configuration at~$q$ attains the
profile with the held word tests of $q$ (admissible since $q$ is attained)
and everything else from~$p$. It also
attains every profile obtained from~$p$ by replacing only its held
word-test set with an admissible set strictly between the two held sets,
so a cover adds several word tests only when no such set is admissible. %
If the supports and held word-test sets agree but the held value-test sets
differ, the configuration at~$q$ attains the profile obtained by adding
to~$p$ one value test held at~$q$ but not at~$p$. The cover forces this
profile to equal~$q$. If only partitions differ, replace $F^p_\ell$ by $F^q_\ell$ at the
last level $\ell$ where they differ. The partitions stay nested: earlier
levels of $p$ are coarser than $F^p_\ell$, hence than $F^q_\ell$, while
later ones agree with~$q$. Moreover, $L'$ meets them all, as no level is
finer than in~$q$. The cases are disjoint, and each
is read from the readout table of Remark~\ref{rem:readback}.
\end{proof}

% compressed 2026-11-24: (e) sentences of the necessity, sufficiency and
% complexity paragraphs merged; (d) the restatement "In Theorem directed
% the verifier first compares the profiles" is cut. Every step of the
% climb is kept; Appendix app:acquisition-proofs cites this climb.
\begin{proof}[Proof of Theorem~\ref{thm:directed}]
For necessity, numerical steps keep the profile and profile steps raise
it, so $p_A\preceq p_B$, and every configuration of a sequence passes the
tests of~$p_A$. Its eigenvalue vector is thus a feasible point of
$\Phi_{p_A}$ (Lemma~\ref{lem:separate}). Profile steps keep $\alpha$, and
numerical steps move it as in the proof of Theorem~\ref{thm:fixed-profile},
which gives the equality of signed cores.

For sufficiency, $\theta_B$ and $\alpha_B$ pass the source's tests (which
$p_B$ holds); since the register conditions are separate, they form with
the source register a feasible configuration~$A'$ at~$p_A$. As
$\alpha_{A'}=\alpha_B$, Theorem~\ref{thm:fixed-profile} joins $A$ to~$A'$.

It remains to climb. Take a saturated chain
$p_A=q_0\prec q_1\prec\dots\prec q_k=p_B$ of attained profiles, each
covering the previous one. It exists because there are finitely many
profiles; it is a chain in the order, not a sequence of steps. At each
$q_j$, pick the register $L_j$ of some feasible configuration (with $L_0$
the source register). Since $q_j\preceq p_B$, every test held at $q_j$ is held
at $p_B$ and passed by $\theta_B$ and $\alpha_B$. Together with $L_j$,
these values form a feasible configuration $A_j$ at~$q_j$, with $A_0=A'$.
If the support shrinks from $q_{j-1}$ to $q_j$, the profile step
$A_{j-1}\to A_j$ keeps the projected tables. If the supports agree, first
interpolate from $L_{j-1}$ to $L_j$ at $q_{j-1}$ by
Lemma~\ref{lem:transfer}(d), as $L_j$ meets the weaker partition
conditions with the same support, then take the profile step with the full
parameter vector fixed. A final interpolation at~$p_B$ replaces $L_k$
by~$L_B$, which also handles $p_A=p_B$.

For the complexity, comparing supports, partitions and test lists checks
$p_A\preceq p_B$ in polynomial time; a failed comparison certifies no, and
otherwise the verifier of Theorem~\ref{thm:certificates}(c) runs on
$\Phi_{p_A}$. For endpoint data supplied as tables, the rational row sums
and divisions have polynomial binary cost, since a common denominator of a
row sum has binary length at most the sum of the input denominator
lengths. At equal profiles the two problems agree
(Remark~\ref{rem:inputs}(ii)), so Theorem~\ref{thm:p-hard} gives the
single-word lower bound with profile steps, and the same transfer from
Theorem~\ref{thm:boundary}(ii), proved in Section~\ref{sec:hard}, gives
the widened one.
\end{proof}

\begin{remark}[What the data determine]\label{rem:readback}
A configuration's \defterm{readout table} lists its answers to three kinds
of \defterm{named queries}: the value of each transition-table entry,
whether each decision is available at each level, and whether each listed
test is held. For a fixed shared-probe model, the readout table determines
the configuration; the probabilities that the held word tests compare need
not. In the sign example, $A^+$ and~$B$ agree on the joint probabilities
with each final base state in the held word test. Their word probabilities
agree because the label probabilities agree, and the projected distributions
agree because the sign-changing eigenvalue enters squared. The
correspondence with attained readout tables preserves and reflects steps,
step types and reconfiguration sequences. %
The same holds under any relabeling that identifies the full attained
profile order, the named queries and their answers, which level is final,
and how transition entries are grouped into each projected sum.
\end{remark}

% compressed 2026-11-24: (c) the row-sum recovery of theta and alpha cites
% the proof of Lemma transfer(a) instead of repeating it; (e) sentences
% merged.
\begin{proof}[Reconstruction in Remark~\ref{rem:readback}]
The positive entries give the support, and two register values share a
block exactly when every available decision agrees on them. From any
starting state, the row sums of the proof of Lemma~\ref{lem:transfer}(a)
give $\theta_m$ and $\theta_m\alpha_{mi}$, hence~$\alpha_{mi}$; then $Q_m$
follows, and dividing each entry by $\theta_mQ_m(g,h)>0$ gives~$L_m$. Held
membership is read from the named answers, so the readout table recovers
the profile and the full parameter vector. Numerical steps are tested on the
recorded entries, and profile steps by attained covers and equality of the
recorded full tables or their projected sums. These checks also determine
the step type and transfer finite sequences in both directions. A relabeling
preserves these checks when it identifies the attained order, the named
queries and their answers, including the final level and the grouping of
entries in each projected sum.
\end{proof}

% compressed 2026-11-24: (e) two sentences merged; (d) the read-back "The
% endpoints still differ in one sign and depend only on the gate count, but
% the instances have no profiles" is cut (Section 4.5 states it).
\paragraph{The algebraic transfer of Section~\ref{sec:circuit-transfer}.}
The emitted system has the $G+2$ eigenvalue coordinates, the probe equations
of the $2G+1$ word tests, the calibration $z=1/4$ and the budgets
$|x_m|\le1/2$; identity reversal adds only identities. Every coordinate
lies in a budget, so the feasible set is bounded; the endpoints are
feasible whatever the circuit, and their signed cores agree exactly when
the output is~$1$, by the proof of Theorem~\ref{thm:p-hard}. Converting each length-three word to an exponent record uses
bounded counts and log-space scans. By Corollary~\ref{cor:real}, with no
boundedness needed, equal signed cores also give the real path-component
verdict, and the endpoints are feasible over the reals.

\subsubsection{Serialization of the reduction of Theorem~\ref{thm:p-hard}}\label{app:serialization}

% compressed 2026-11-24: (e) merges, and the record items start with their
% names instead of "A header gives" etc.; the record layout, field
% bounds, work-tape analysis and the malformed-input handling are kept in
% full.
Both parts of Theorem~\ref{thm:boundary} read one explicit binary grammar,
in which a word test lists one word on one side and one or two on the
other. Goldschlager's circuits come with an assignment to their input
variables~\cite{goldschlager}; substituting the assigned constants needs
one scan of the assignment and an index counter. The transducer then
writes the following records, with binary natural fields separated by
delimiters that fixed binary codes can replace.
\begin{itemize}
\item Header: $G+2$ labels, two base states, one register value, one
probe, $H=3$, $\delta=1/(G+2)$, $\kappa=1/2$, $2G+1$ word tests and one
calibration.
\item Geometry: the probe's signs, the register weight, the single level
and its partition, and the single admissible set.
\item Per label: its number, its reversal, its probability, its eigenvalue
at each endpoint, its four unit register entries and its four support
flags.
\item Per word test: its number, its two words as triples of label numbers,
and three membership flags, all one: in the admissible set and in the held
tests of each endpoint. Calibration: the anchor, the probe, the value
$1/4$ and two held flags.
\end{itemize}
The flags encode the explicitly listed admissible set and endpoint held
sets, which rescanning can also emit as lists. Records appear in order:
header, geometry, labels, output test, the two slots of each gate in source
order, and calibration. Gate-label numbers and gate-slot test numbers
are fixed affine functions of the gate index, with a fixed offset for each
slot. The anchor, sentinel and output test have fixed numbers. There are $3G+6$
records, each with at most seventeen fields; the $2G+16$ field bound
applies to the parameter records. Entries shared by the endpoints are
printed twice where the grammar asks for separate tables. Requested
transition-table entries are $3/(8(G+2))$ or $5/(8(G+2))$, by the probe
signs and the sentinel's sign, so a denominator $8(G+2)$ still has
logarithmic binary length.

For an input of length~$\ell_{\mathrm{in}}$, the work tape holds a constant
number of counters and fields of binary length
$O(\log(\ell_{\mathrm{in}}+2))$. They record the input position, the
current record and slot, the header data, the current gate record and the
current output field and digit. A slot is emitted by scanning to the requested gate
record, reading its type and at most two input indices, and computing
label and test numbers by addition and doubling. Each field is written in
binary from its leading digit, with a digit counter. The output has binary
length $O((G+1)\log(G+2))$, and the running time is polynomial.

First count the input length $\ell_{\mathrm{in}}$ and scan the gate
records. Check that there are $G\ge1$ records, that each has an allowed tag
and exactly the required fields, that the output index is in
$1,\ldots,G$, and that every predecessor is in $1,\ldots,i-1$ at
record~$i$. Reject an index as soon as its value exceeds
$\ell_{\mathrm{in}}$. On failure emit the instance for the one-gate
constant-$0$ circuit, a fixed feasible negative instance. So the reduction
is total, and its scans use logarithmic space on every input.

